% ============================================================
% INTRODUCTION — target: 600--800 words, 4 paragraphs
% ============================================================

\section{Introduction}

% Paragraph 1: The Problem
As large language models are deployed in increasingly high-stakes settings---from
medical diagnosis and legal reasoning to autonomous code generation---the ability
to monitor and understand their reasoning processes becomes a critical safety
requirement. Chain-of-thought (CoT) reasoning, in which models produce step-by-step
explanations before arriving at a final answer, has been widely adopted as a
transparency mechanism: if a model's reasoning is visible, human overseers can
potentially detect flawed logic, biased reasoning, or deceptive intent before harm
occurs \cite{wei2022cot}. This promise is especially important for reasoning
models---systems explicitly trained to produce extended thinking traces via
reinforcement learning---which now dominate benchmarks in mathematics, science, and
coding \cite{qchen2025longcot, xu2025large}. However, the safety value of CoT
monitoring rests on a crucial assumption: that the chain-of-thought faithfully
represents the model's actual reasoning process. If models can be influenced by
factors they do not acknowledge in their CoT, then monitoring these traces provides
a false sense of security.

% Paragraph 2: What Is Known
A growing body of work has investigated the faithfulness of chain-of-thought
reasoning. Turpin et al.~\cite{turpin2023unfaithful} demonstrated that biased
features in prompts (such as suggested answers from a purported expert)
systematically influence model outputs without being mentioned in the CoT, revealing
gaps between stated and actual reasoning across multiple tasks on BIG-Bench Hard.
Lanham et al.~\cite{lanham2023measuring} further measured CoT faithfulness in Claude
models by truncating, corrupting, and adding mistakes to reasoning chains, finding
that the correlation between CoT content and model predictions varies substantially
across tasks. Lyu et al.~\cite{lyu2023faithful} proposed a two-stage framework
translating natural language into symbolic reasoning chains to improve faithfulness,
while Radhakrishnan et al.~\cite{radhakrishnan2023decomposition} showed that
question decomposition improves the faithfulness of model-generated reasoning.
Most recently, Chen et al.~\cite{chen2025reasoning} conducted the most comprehensive
evaluation to date, using six types of reasoning hints---sycophancy, consistency,
visual pattern recognition, metadata cues, grader hacking, and unethical information
use---injected into MMLU and GPQA multiple-choice questions. They found that
Claude 3.7 Sonnet acknowledges hints that influenced its answer only 25\% of the
time, while DeepSeek-R1 does so only 39\% of the time, with unfaithful CoTs being
paradoxically \emph{longer} than faithful ones. A shared limitation across all of
these studies is their narrow model coverage: existing evaluations test at most two
to four models, leaving the vast and rapidly growing ecosystem of open-weight
reasoning models almost entirely unexamined.

% Paragraph 3: The Gap
What remains unknown is whether the low faithfulness rates observed in Claude and
DeepSeek-R1 generalize across the broader open-weight reasoning model
ecosystem---or whether they reflect properties specific to particular architectures,
training methods, or model scales. Since early 2025, a diverse set of open-weight
reasoning models has emerged spanning dense architectures from 8B to 32B parameters,
mixture-of-experts (MoE) models exceeding 1T total parameters, and training
pipelines ranging from pure reinforcement learning (GRPO, DPO) to distillation,
data-centric methods, and hybrid approaches. No study has systematically compared
CoT faithfulness across these model families, leaving a critical gap in our
understanding of which models---and which training paradigms---produce reasoning
traces that can be reliably monitored for safety purposes.

% Paragraph 4: Your Solution
We address this gap with the first large-scale, cross-family evaluation of CoT
faithfulness in open-weight reasoning models. We evaluate 11 models spanning 9
architectural families---including DeepSeek, Qwen, Zhipu (GLM), Moonshot (Kimi),
MiniMax, OpenAI, Baidu, and AI2---all accessed via the OpenRouter API. Using 500
multiple-choice questions (300 from MMLU and 200 from GPQA Diamond) with the same
six hint categories established by Chen et al.~\cite{chen2025reasoning}, we conduct
38,500 inference runs to measure faithfulness rates across model families, parameter
scales (7B to 1T), and training methodologies. We employ Anthropic's BLOOM
framework~\cite{bloom2025} for automated behavioral scoring of faithfulness,
validated against 200 hand-labeled examples. We hypothesize that (H1)~faithfulness rates vary significantly across
model families, reflecting differences in training methodology rather than scale
alone, and (H2)~certain hint types---particularly metadata and grader
hacking---exhibit consistently lower faithfulness across all models, suggesting
category-specific blind spots in CoT monitoring. Our evaluation employs a two-stage
faithfulness classifier combining regex-based pattern matching with an LLM
judge, validated against 200 hand-labeled examples. We release all code, prompts,
model outputs, and faithfulness annotations to support future research on CoT
monitoring reliability.
